% Options for packages loaded elsewhere
\PassOptionsToPackage{unicode}{hyperref}
\PassOptionsToPackage{hyphens}{url}
\PassOptionsToPackage{table}{xcolor}
%
\documentclass[
  letterpaper,
]{article}
\usepackage{amsmath,amssymb}
\usepackage{iftex}
\usepackage[T1]{fontenc}
\usepackage[utf8]{inputenc}
\usepackage{textcomp}
\usepackage{lmodern}
\ifPDFTeX\else
  % xetex/luatex font selection
\fi
% Use upquote if available, for straight quotes in verbatim environments
\IfFileExists{upquote.sty}{\usepackage{upquote}}{}
\IfFileExists{microtype.sty}{% use microtype if available
  \usepackage[]{microtype}
  \UseMicrotypeSet[protrusion]{basicmath} % disable protrusion for tt fonts
}{}
\makeatletter
\@ifundefined{KOMAClassName}{% if non-KOMA class
  \IfFileExists{parskip.sty}{%
    \usepackage{parskip}
  }{% else
    \setlength{\parindent}{0pt}
    \setlength{\parskip}{6pt plus 2pt minus 1pt}}
}{% if KOMA class
  \KOMAoptions{parskip=half}}
\makeatother
\usepackage{xcolor}
\usepackage{color}
\usepackage{fancyvrb}
\newcommand{\VerbBar}{|}
\newcommand{\VERB}{\Verb[commandchars=\\\{\}]}
\DefineVerbatimEnvironment{Highlighting}{Verbatim}{commandchars=\\\{\}}
% Add ',fontsize=\small' for more characters per line
\newenvironment{Shaded}{}{}
\newcommand{\AlertTok}[1]{\textcolor[rgb]{1.00,0.00,0.00}{\textbf{#1}}}
\newcommand{\AnnotationTok}[1]{\textcolor[rgb]{0.38,0.63,0.69}{\textbf{\textit{#1}}}}
\newcommand{\AttributeTok}[1]{\textcolor[rgb]{0.49,0.56,0.16}{#1}}
\newcommand{\BaseNTok}[1]{\textcolor[rgb]{0.25,0.63,0.44}{#1}}
\newcommand{\BuiltInTok}[1]{\textcolor[rgb]{0.00,0.50,0.00}{#1}}
\newcommand{\CharTok}[1]{\textcolor[rgb]{0.25,0.44,0.63}{#1}}
\newcommand{\CommentTok}[1]{\textcolor[rgb]{0.38,0.63,0.69}{\textit{#1}}}
\newcommand{\CommentVarTok}[1]{\textcolor[rgb]{0.38,0.63,0.69}{\textbf{\textit{#1}}}}
\newcommand{\ConstantTok}[1]{\textcolor[rgb]{0.53,0.00,0.00}{#1}}
\newcommand{\ControlFlowTok}[1]{\textcolor[rgb]{0.00,0.44,0.13}{\textbf{#1}}}
\newcommand{\DataTypeTok}[1]{\textcolor[rgb]{0.56,0.13,0.00}{#1}}
\newcommand{\DecValTok}[1]{\textcolor[rgb]{0.25,0.63,0.44}{#1}}
\newcommand{\DocumentationTok}[1]{\textcolor[rgb]{0.73,0.13,0.13}{\textit{#1}}}
\newcommand{\ErrorTok}[1]{\textcolor[rgb]{1.00,0.00,0.00}{\textbf{#1}}}
\newcommand{\ExtensionTok}[1]{#1}
\newcommand{\FloatTok}[1]{\textcolor[rgb]{0.25,0.63,0.44}{#1}}
\newcommand{\FunctionTok}[1]{\textcolor[rgb]{0.02,0.16,0.49}{#1}}
\newcommand{\ImportTok}[1]{\textcolor[rgb]{0.00,0.50,0.00}{\textbf{#1}}}
\newcommand{\InformationTok}[1]{\textcolor[rgb]{0.38,0.63,0.69}{\textbf{\textit{#1}}}}
\newcommand{\KeywordTok}[1]{\textcolor[rgb]{0.00,0.44,0.13}{\textbf{#1}}}
\newcommand{\NormalTok}[1]{#1}
\newcommand{\OperatorTok}[1]{\textcolor[rgb]{0.40,0.40,0.40}{#1}}
\newcommand{\OtherTok}[1]{\textcolor[rgb]{0.00,0.44,0.13}{#1}}
\newcommand{\PreprocessorTok}[1]{\textcolor[rgb]{0.74,0.48,0.00}{#1}}
\newcommand{\RegionMarkerTok}[1]{#1}
\newcommand{\SpecialCharTok}[1]{\textcolor[rgb]{0.25,0.44,0.63}{#1}}
\newcommand{\SpecialStringTok}[1]{\textcolor[rgb]{0.73,0.40,0.53}{#1}}
\newcommand{\StringTok}[1]{\textcolor[rgb]{0.25,0.44,0.63}{#1}}
\newcommand{\VariableTok}[1]{\textcolor[rgb]{0.10,0.09,0.49}{#1}}
\newcommand{\VerbatimStringTok}[1]{\textcolor[rgb]{0.25,0.44,0.63}{#1}}
\newcommand{\WarningTok}[1]{\textcolor[rgb]{0.38,0.63,0.69}{\textbf{\textit{#1}}}}
\usepackage{longtable,booktabs,array}
\usepackage{calc} % for calculating minipage widths
% Correct order of tables after \paragraph or \subparagraph
\usepackage{etoolbox}
\makeatletter
\patchcmd\longtable{\par}{\if@noskipsec\mbox{}\fi\par}{}{}
\makeatother
% Allow footnotes in longtable head/foot
\IfFileExists{footnotehyper.sty}{\usepackage{footnotehyper}}{\usepackage{footnote}}
\makesavenoteenv{longtable}
\setlength{\emergencystretch}{3em} % prevent overfull lines
\providecommand{\tightlist}{%
  \setlength{\itemsep}{0pt}\setlength{\parskip}{0pt}}
\setcounter{secnumdepth}{5}
% Self-contained presentation rules embedded in every generated LaTeX file.
\usepackage{iftex}
\ifPDFTeX\else\PackageError{SSE-ASAS}{Compile with pdfLaTeX only}{}\fi
\usepackage[T1]{fontenc}
\usepackage[utf8]{inputenc}
\usepackage{lmodern}
\usepackage{microtype}
\usepackage{geometry}
\geometry{margin=0.93in,headheight=16pt}
\usepackage{xcolor} % [table] passed at top of file (option-clash fix)
\definecolor{ink}{HTML}{15283B}
\definecolor{brand}{HTML}{075B75}
\definecolor{brandlight}{HTML}{EAF5F8}
\definecolor{warm}{HTML}{A15C1D}
\definecolor{muted}{HTML}{536577}
\usepackage{fancyhdr}
\usepackage{titlesec}
\usepackage{enumitem}
\usepackage{booktabs}
\usepackage{longtable}
\usepackage{array}
\usepackage{tabularx}
\usepackage{tcolorbox}
\tcbuselibrary{breakable,skins}
\usepackage{listings}
\usepackage{fancyvrb}
\usepackage{hyperref}
\hypersetup{colorlinks=true,linkcolor=brand,urlcolor=brand,citecolor=brand}
\urlstyle{same}
\usepackage{lastpage}
\setlength{\emergencystretch}{3em}
\setlength{\parskip}{0.45em}
\setlength{\parindent}{0pt}
\setlist[itemize]{leftmargin=*,topsep=3pt,itemsep=2pt}
\setlist[enumerate]{leftmargin=*,topsep=3pt,itemsep=2pt}
\titleformat{\section}{\Large\bfseries\color{ink}}{\thesection}{0.75em}{}[\vspace{-2pt}\color{brand}\titlerule]
\titleformat{\subsection}{\large\bfseries\color{brand}}{\thesubsection}{0.7em}{}
\titleformat{\subsubsection}{\normalsize\bfseries\color{ink}}{\thesubsubsection}{0.7em}{}
\pagestyle{fancy}
\fancyhf{}
\fancyhead[L]{\small\sffamily\textcolor{brand}{SOFTWARE SECURITY ENGINEERING}}
\fancyhead[R]{\small\sffamily\textcolor{muted}{TECHNICAL SERIES}}
\fancyfoot[L]{\scriptsize\sffamily Version 1.0 \textbullet\ 14 October 2026}
\fancyfoot[R]{\scriptsize\sffamily Page \thepage\ of \pageref{LastPage}}
\renewcommand{\headrulewidth}{0.5pt}
\renewcommand{\footrulewidth}{0.2pt}
\setcounter{tocdepth}{2}
\setcounter{secnumdepth}{3}
\lstdefinestyle{sse}{basicstyle=\ttfamily\footnotesize,frame=single,framerule=0.4pt,rulecolor=\color{brand!40},backgroundcolor=\color{brandlight!45},numbers=left,numberstyle=\scriptsize\color{muted},numbersep=6pt,breaklines=true,breakatwhitespace=false,tabsize=2,showstringspaces=false,columns=fullflexible,keepspaces=true,captionpos=b,aboveskip=9pt,belowskip=9pt}
\lstset{style=sse}
\newtcolorbox{sourceconcept}[1][]{enhanced,breakable,colback=brandlight,colframe=brand!65!black,title={ORIGINAL CONCEPT},fonttitle=\bfseries\sffamily,#1}
\newtcolorbox{modernextension}[1][]{enhanced,breakable,colback=white,colframe=warm!85!black,title={MODERN EXTENSION},fonttitle=\bfseries\sffamily,#1}
\newtcolorbox{fieldwork}[1][]{enhanced,breakable,colback=white,colframe=ink!55,title={PRACTICAL APPLICATION},fonttitle=\bfseries\sffamily,#1}
\AtBeginDocument{\color{ink}}
\makeatletter
\renewcommand{\maketitle}{%
 \begin{titlepage}\thispagestyle{empty}
 \vspace*{1.5cm}{\color{brand}\rule{\textwidth}{3pt}}\par\vspace{1cm}
 {\sffamily\bfseries\color{muted}\large SOFTWARE SECURITY ENGINEERING\\ APPLICATION SECURITY ASSESSMENT SERIES}\par\vspace{1.05cm}
 {\sffamily\bfseries\color{ink}\fontsize{27}{34}\selectfont\@title\par}\vspace{0.8cm}
 {\sffamily\large\color{brand}\@author\par}\vfill
 \begin{tcolorbox}[colback=brandlight,colframe=brand!55,boxrule=0.5pt]
 \small\textbf{Editorial policy}\quad Original teaching synthesis with book-grounded concepts distinguished from contemporary practices. Modern platform behavior is version-sensitive; validate against linked vendor documentation before production use.
 \end{tcolorbox}
 {\color{brand}\rule{\textwidth}{2pt}}\par\vspace{0.3cm}
 {\sffamily\small Version 1.0 \quad|\quad October 2026 \quad|\quad pdfLaTeX}\end{titlepage}\setcounter{page}{1}}
\makeatother
\ifLuaTeX
  \usepackage{selnolig}  % disable illegal ligatures
\fi
\usepackage{bookmark}
\IfFileExists{xurl.sty}{\usepackage{xurl}}{} % add URL line breaks if available
\urlstyle{same}
\hypersetup{
  pdftitle={Application Security Assessment Procedures},
  pdfauthor={Independent Technical Study Edition},
  hidelinks,
  pdfcreator={LaTeX via pandoc}}

\title{Application Security Assessment Procedures}
\author{Independent Technical Study Edition}
\date{}

\begin{document}
\maketitle

{
\setcounter{tocdepth}{3}
\tableofcontents
}
\section{Standard operating procedure
(SOP)}\label{standard-operating-procedure-sop}

This guide defines a reproducible, authorization-aware application
security assessment from intake through closure. It combines manual
review and program analysis under a shared evidence model. It is a
\textbf{methodology}, not a claim of certification or regulatory
compliance; confirm your organization's retention, privacy, and
incident-response requirements before use.

\textbf{Book foundation:} Dowd et al., Chapters 2--4; Bishop, Chapters
19--25; Chess and West, Chapters 3--4. \textbf{Modern extension:} OWASP
ASVS 5.0.0, SSDF 1.1, GitHub Code Security, SARIF, and risk-based
remediation governance.

\section{Process topology and control
gates}\label{process-topology-and-control-gates}

\begin{enumerate}
\def\labelenumi{\arabic{enumi}.}
\tightlist
\item
  \textbf{Assessment authorized}: signed/recorded scope, systems,
  branches, environments, test limitations, privacy rules.
\item
  \textbf{Assets and invariants recorded}: architecture, ownership,
  sensitive data, threat model, trust boundaries.
\item
  \textbf{Evidence collection enabled}: repo snapshot, build
  instructions, test account, scan entitlement, log handling.
\item
  \textbf{Static baseline established}: verify extractors, language
  coverage, query suite, baseline findings.
\item
  \textbf{Manual review performed}: priority paths, business rules,
  auth, native/API/platform review.
\item
  \textbf{Findings verified}: trace, repeatability, impact, mapping,
  owner and confidence.
\item
  \textbf{Remediation planned}: preferred design fix, compensating
  controls, dates, and exception process.
\item
  \textbf{Fix validated}: regression tests, rescan of fix branch or PR,
  review of adjacent paths.
\item
  \textbf{Assessment accepted}: decision on residual findings by
  authorized owner; report archived under retention rules.
\end{enumerate}

The phrase ``fix branch validated'' is intentional: a clean default
branch after merge does not prove the fix was scanned or required
\textbf{before} merging.

\section{Phase I: intake, scope, and
architecture}\label{phase-i-intake-scope-and-architecture}

\subsection{Engagement charter
template}\label{engagement-charter-template}

\begin{longtable}[]{@{}
  >{\raggedright\arraybackslash}p{(\columnwidth - 2\tabcolsep) * \real{0.5000}}
  >{\raggedright\arraybackslash}p{(\columnwidth - 2\tabcolsep) * \real{0.5000}}@{}}
\toprule\noalign{}
\begin{minipage}[b]{\linewidth}\raggedright
Field
\end{minipage} & \begin{minipage}[b]{\linewidth}\raggedright
Required entry
\end{minipage} \\
\midrule\noalign{}
\endhead
\bottomrule\noalign{}
\endlastfoot
Application / business owner & Name, team, contact and escalation \\
Repositories / SHAs & Exact commit IDs, default and protected
branches \\
Components & APIs, workers, front end, dependencies, infra \\
Environments & Dev, staging, pre-production; exclude production unless
approved \\
Authorized methods & Static review, SAST, authorized functional
validation \\
Prohibited methods & DoS, lateral movement, secret exfiltration,
unapproved fuzzing \\
Data safeguards & Data classification, evidence redaction, retention \\
Decision authority & Who accepts findings, exceptions, and closure \\
\end{longtable}

\subsection{Threat-model worksheet}\label{threat-model-worksheet}

For each data flow, capture
\texttt{\{source,\ principal,\ trust\ boundary,\ destination,\ sensitive\ asset,\ expected\ decision,\ existing\ control,\ assumption,\ test\}}.
Prioritize money movement, cross-tenant or cross-user object access,
admin control, credential use, external calls, parsers, and persistent
state transitions. Include asynchronous paths and direct data-store
interactions.

\section{Phase II: coverage planning}\label{phase-ii-coverage-planning}

\textbf{Test planning:} select one or more assurance requirements per
security-sensitive entry point. For Web/API systems use version-pinned
ASVS requirements where applicable. Use CWE IDs for weakness
classification. Coverage must include manual business-logic checks
because static tools cannot establish every authorization or workflow
invariant.

\begin{longtable}[]{@{}
  >{\raggedright\arraybackslash}p{(\columnwidth - 6\tabcolsep) * \real{0.2500}}
  >{\raggedright\arraybackslash}p{(\columnwidth - 6\tabcolsep) * \real{0.2500}}
  >{\raggedright\arraybackslash}p{(\columnwidth - 6\tabcolsep) * \real{0.2500}}
  >{\raggedright\arraybackslash}p{(\columnwidth - 6\tabcolsep) * \real{0.2500}}@{}}
\toprule\noalign{}
\begin{minipage}[b]{\linewidth}\raggedright
Asset / workflow
\end{minipage} & \begin{minipage}[b]{\linewidth}\raggedright
Manual strategy
\end{minipage} & \begin{minipage}[b]{\linewidth}\raggedright
Tool evidence
\end{minipage} & \begin{minipage}[b]{\linewidth}\raggedright
Negative test
\end{minipage} \\
\midrule\noalign{}
\endhead
\bottomrule\noalign{}
\endlastfoot
Tenant invoice lookup & Trace principal-to-query & CodeQL query alert
review & Cross-tenant read denied \\
Payment approval & State-machine review & Dependency and code-scan
results & Submitter cannot approve \\
File upload & Parser and path audit & Native/tool findings & Traversal
and oversized file rejected \\
Worker credential & IAM / secret review & Secret scanning inventory &
Unauthorized scope not assumed \\
External webhook & Origin, idempotency, replay checks & Code scan plus
integration test & Replay does not duplicate settlement \\
\end{longtable}

\subsection{SAST baseline and blind
spots}\label{sast-baseline-and-blind-spots}

Record the scanning configuration and \textbf{actual} language and file
coverage. Verify action execution, permission scopes, branches, build
mode, extraction failures, excluded paths, custom rule pack versions,
and success criteria. Differentiate alerts open on default branch from
PR alerts. Evaluate a representative seeded issue in an isolated test
repository to confirm the alert and required check behave as intended.

\section{Phase III: investigation and finding
confirmation}\label{phase-iii-investigation-and-finding-confirmation}

\begin{enumerate}
\def\labelenumi{\arabic{enumi}.}
\tightlist
\item
  Review trace from untrusted source or unauthorized subject to
  protected operation.
\item
  Distinguish system preconditions from attack payload and separate
  hypothesis from observation.
\item
  Reproduce only under authorized safe conditions; do not perform
  destructive proof-of-concept tests.
\item
  Determine whether an existing control actually blocks the path at the
  relevant layer.
\item
  Map to precise CWE and versioned verification requirement; record
  disagreements.
\item
  Identify root cause and adjacent routes sharing the same pattern.
\item
  Document a regression test and remediation strategy before a finding
  is closed.
\end{enumerate}

\subsection{Finding record template}\label{finding-record-template}

\textbf{ID:} APPSEC-YYYY-NNN. \textbf{Title:} component plus failure
mechanism. \textbf{Source:} immutable repo/commit/line and scanner rule
if any. \textbf{Affected versions:} evidence-based, not guessed.
\textbf{Invariant:} expected security property and authority.
\textbf{Trace:} actor \(\to\) entry \(\to\) decision \(\to\) sensitive
operation. \textbf{Reproduction:} authorized conditions and redacted
result. \textbf{Impact:} assets and harm. \textbf{Likelihood:} exposure,
required conditions and uncertainty. \textbf{References:} CWE, ASVS ID
(versioned), code scan alert URL. \textbf{Resolution:} code/control,
owner, due date, test, re-scan evidence. \textbf{Residual risk:}
acknowledged by decision authority.

\begin{fieldwork}
**Sample finding:** "Invoice export omits tenant scoping." An authenticated tenant-A user can export a tenant-B invoice if supplied its identifier. Expected: tenant context is applied at the persistent query boundary. Root cause: export worker uses a privileged fetch-by-ID path. Validation: integration test fails before patch; after patch, unauthorized export returns a non-disclosing response and the job records denial. Check similar print/email export paths.
\end{fieldwork}

\section{Phase IV: disposition and risk
decisions}\label{phase-iv-disposition-and-risk-decisions}

Classification is not a substitute for fixing a defect. Define a policy
for states such as open, confirmed, remediated, verified, accepted risk,
duplicate, invalid finding, and test-only code. The GitHub UI's labels
and capabilities are platform-specific; organizational approvals and
evidence live in your system of record.

\begin{longtable}[]{@{}
  >{\raggedright\arraybackslash}p{(\columnwidth - 2\tabcolsep) * \real{0.5000}}
  >{\raggedright\arraybackslash}p{(\columnwidth - 2\tabcolsep) * \real{0.5000}}@{}}
\toprule\noalign{}
\begin{minipage}[b]{\linewidth}\raggedright
Decision
\end{minipage} & \begin{minipage}[b]{\linewidth}\raggedright
Evidence required
\end{minipage} \\
\midrule\noalign{}
\endhead
\bottomrule\noalign{}
\endlastfoot
Confirmed & Reachable flaw or demonstrable violation; trace and
impact \\
False positive & Tool model invalid for the actual code path; show
proof \\
Mitigated & Concrete control interrupts the path; verify control and
residual risk \\
Used in tests & Clearly isolated test-only path cannot ship/expose
assets \\
Accepted risk & Named authority, business rationale, expiration/review
date \\
Fixed & Patch merged under required controls; tests and appropriate scan
evidence \\
\end{longtable}

Avoid using ``won't fix'' as a generic substitute for a documented risk
acceptance; map any platform-specific dismissal options to the
organization's approved workflow.

\subsection{Prioritization}\label{prioritization}

Evaluate impact, exploit preconditions, asset criticality, reachability,
exposure, and mitigation. Severity from a tool, CVSS assessment, and
exploit probability are different signals. Do not compute spurious
decimal precision from subjective estimates. Establish an escalation
rule for live credential exposure, active compromise, or high-impact
unauthorized data access.

\section{Phase V: remediation verification and
closure}\label{phase-v-remediation-verification-and-closure}

A fix is verified when the identified cause is removed or meaningfully
constrained, authorized negative tests pass, relevant
dependency/build/scan checks complete on the correct ref, and adjacent
endpoints are checked. Scanners can report ``fixed'' because a path
changed or disappeared from extracted source; inspect the actual
security invariant as well.

\textbf{Release gate:} protected branch requires successful build/tests
and appropriate code scanning check. Ensure scan workflow triggers PRs
and does not skip due to filters or permissions. For manual-build
languages, verify the files of interest are included in the CodeQL
database; for third-party tools validate SARIF upload and run category
association.

\textbf{Acceptance packet:} scope version; coverage summary; findings
ledger; evidence links; exception decisions; fix verification;
unresolved assumptions; owner sign-off. Mark unknowns explicitly; never
overstate assurance from a finite assessment.

\section{Evidence quality metrics and
governance}\label{evidence-quality-metrics-and-governance}

\begin{longtable}[]{@{}
  >{\raggedright\arraybackslash}p{(\columnwidth - 4\tabcolsep) * \real{0.3333}}
  >{\raggedright\arraybackslash}p{(\columnwidth - 4\tabcolsep) * \real{0.3333}}
  >{\raggedright\arraybackslash}p{(\columnwidth - 4\tabcolsep) * \real{0.3333}}@{}}
\toprule\noalign{}
\begin{minipage}[b]{\linewidth}\raggedright
Metric
\end{minipage} & \begin{minipage}[b]{\linewidth}\raggedright
Definition
\end{minipage} & \begin{minipage}[b]{\linewidth}\raggedright
Failure mode to avoid
\end{minipage} \\
\midrule\noalign{}
\endhead
\bottomrule\noalign{}
\endlastfoot
Assessment coverage & Tested controls / in-scope applicable controls &
Treating untested as passed \\
Triage latency & Detection to first justified classification &
Optimizing by dismissing rapidly \\
Remediation latency & Confirmed to verified fix & Ignoring risk while
awaiting release \\
Scan health & Successful analysis / scheduled scans & Counting no-op
workflows as success \\
Reopen rate & Previously closed flaws that recur & Blaming teams instead
of root causes \\
Exception aging & Active exceptions past review date & Permanent
invisible accepted risk \\
\end{longtable}

\section{Ready-to-use artifacts}\label{ready-to-use-artifacts}

\textbf{Assessment sign-off checklist:} scope authorized; commit
snapshot pinned; owners recorded; architecture and trust boundaries
reviewed; SAST successful and coverage measured; manual hypotheses
executed; findings independently traceable; remediation plans contain
regressions; secrets handled safely; exceptions signed and dated; report
indexed; retention and follow-up dates assigned.

\textbf{Peer review quality gate:} another engineer can reproduce the
reasoning from the finding packet without privileged access to the first
reviewer's workstation. If the result depends on an undisclosed
assumption, mark confidence lower and request clarification.

\section{Annex A: assessment artifacts to fill
in}\label{annex-a-assessment-artifacts-to-fill-in}

\subsection{Security assessment plan
(SAP)}\label{security-assessment-plan-sap}

\textbf{System boundary:} repositories, running services, deployment
dependencies and trust relationships. \textbf{Business objectives:} key
transactions and acceptable residual risk. \textbf{Assessment
objective:} named security properties and versioned ASVS requirements.
\textbf{In-scope evidence:} code, tests, configuration, build artifacts
and telemetry available under authorization. \textbf{Methods:} design
review, code trace, static queries, controlled negative testing.
\textbf{Team:} assessor, peer reviewer, development owner, risk owner.
\textbf{Milestones:} scope freeze, preliminary findings, remediation
review, validation, acceptance. \textbf{Exclusions:} out-of-scope
assets, prohibited testing, third-party components without rights,
production systems if not approved.

A concrete plan should record the selected commit SHA and environment
state; otherwise two engineers may inspect different code while
describing the same review. When the build depends on external services,
capture dependency versions and distinguish attempted versus completed
test coverage.

\subsection{Requirements-to-test traceability
matrix}\label{requirements-to-test-traceability-matrix}

\begin{longtable}[]{@{}
  >{\raggedright\arraybackslash}p{(\columnwidth - 8\tabcolsep) * \real{0.2000}}
  >{\raggedright\arraybackslash}p{(\columnwidth - 8\tabcolsep) * \real{0.2000}}
  >{\raggedright\arraybackslash}p{(\columnwidth - 8\tabcolsep) * \real{0.2000}}
  >{\raggedright\arraybackslash}p{(\columnwidth - 8\tabcolsep) * \real{0.2000}}
  >{\raggedright\arraybackslash}p{(\columnwidth - 8\tabcolsep) * \real{0.2000}}@{}}
\toprule\noalign{}
\begin{minipage}[b]{\linewidth}\raggedright
Requirement ID
\end{minipage} & \begin{minipage}[b]{\linewidth}\raggedright
Invariant
\end{minipage} & \begin{minipage}[b]{\linewidth}\raggedright
Test technique
\end{minipage} & \begin{minipage}[b]{\linewidth}\raggedright
Evidence
\end{minipage} & \begin{minipage}[b]{\linewidth}\raggedright
Status
\end{minipage} \\
\midrule\noalign{}
\endhead
\bottomrule\noalign{}
\endlastfoot
REQ-AUTH-01 & Every invoice read enforces tenant & Negative integration
& Test run / route trace & Planned \\
REQ-INTEG-02 & Submitter cannot approve & State transition review & Test
and transaction check & Planned \\
REQ-INPUT-03 & Data cannot alter SQL grammar & Static trace and unit
test & Parameter binding and tests & Planned \\
REQ-SECRET-04 & No live token enters source & CI secret scanning & Scan
log and handling SOP & Planned \\
REQ-CI-05 & Required scan executes on PR SHA & Branch protection
verification & Check run and ref & Planned \\
\end{longtable}

Use organization-specific IDs for locally authored requirements. ASVS
control IDs, when included, must retain a version prefix; do not invent
numeric ASVS identifiers. A control that is not applicable must have a
written rationale, not be silently omitted from denominator
calculations.

\subsection{Risk and remediation
ledger}\label{risk-and-remediation-ledger}

\begin{Shaded}
\begin{Highlighting}[]
\NormalTok{Finding ID:          ...}
\NormalTok{Asset and owner:     ...}
\NormalTok{Severity rationale:  asset + impact + preconditions + exposure}
\NormalTok{Exploit confidence:  observed / likely / theoretical}
\NormalTok{CWE mapping:         validated specific ID and source link}
\NormalTok{Detection origin:    manual / CodeQL / other SAST / DAST / IAST}
\NormalTok{Tool rule ID:        ... (if applicable)}
\NormalTok{Fix recommendation:  root{-}cause control; avoid local symptom{-}only patch}
\NormalTok{Assigned engineer:   ...}
\NormalTok{Due date:            ...}
\NormalTok{Compensating control:... (if verified)}
\NormalTok{Reproduction status: ...}
\NormalTok{Test evidence:       ...}
\NormalTok{PR/check run evidence:...}
\NormalTok{Exception authority: ... (if risk accepted)}
\NormalTok{Exception review:    ...}
\NormalTok{Closure reviewer:    ...}
\end{Highlighting}
\end{Shaded}

\subsection{Final report layout}\label{final-report-layout}

\textbf{Executive summary (one page):} application criticality,
assessment period, tested scope, key risks, assurance limits, decisions
needed, and recommended near-term actions. \textbf{Technical summary:}
attack surface, architecture and trust boundaries, methods and actual
test/scan coverage, notable controls, confirmed findings by root-cause
family, systemic recommendations, and unresolved questions.
\textbf{Detailed findings:} trace, repro in approved environment,
business consequences, classification, fix, verification, owner.
\textbf{Appendix:} scope letter, artifacts inventory, test environment,
tools and versions, immutable revisions, exception ledger, evidence
retention policy.

\section{Annex B: release-gate decision
examples}\label{annex-b-release-gate-decision-examples}

\textbf{Case 1 -- high-confidence CodeQL alert on PR.} Developer updates
code; required CodeQL analysis completes on updated PR ref; negative
unit test demonstrates control; security review confirms path no longer
unsafe; branch protection prevents merge until required checks succeed.
On default branch, confirm the resulting revision was analyzed, but do
not treat that as the sole premerge control.

\textbf{Case 2 -- critical design flaw with no scanner alert.} An order
export lacks tenant filtering; no tool identifies the missing invariant.
Manual auditor files a documented finding with negative test. Release
gate depends on the organization's risk acceptance process; the absence
of an alert does not allow auto-approval.

\textbf{Case 3 -- scan completed but zero files analyzed.} Although the
workflow check reports success, extraction coverage does not include the
new language. Treat as a scanning failure/coverage gap, require
investigation and re-analysis, and record it separately from
vulnerability disposition.

\textbf{Case 4 -- unresolved credential alert.} If the credential is
real and reachable, invoke incident handling and revoke/rotate as
appropriate; control cannot be satisfied merely by deleting the latest
line. Validate current repository state, history exposure, downstream
scope, and prevention policy.

\section{Annex C: independent evidence review
checklist}\label{annex-c-independent-evidence-review-checklist}

A second reviewer should ask: Did we have authority to test? Is the
scope stable? Are customer data and secrets handled correctly? Can each
finding be verified on the stated version? Was scanner/tool coverage
genuinely measured? Does each remediation preserve intended operations
while rejecting unwanted ones? Was the risk owner authorized? Were old
issues accidentally re-labeled as fixed because a scan stopped working?
Do findings and metrics identify untested components? Can the
development team use the report without asking the original auditor to
restate every detail?

An evidence record that satisfies these questions creates a reusable
basis for engineering decisions, risk reporting, and future regression
review.

\section{Source and modern extension
map}\label{source-and-modern-extension-map}

\textbf{Book-grounded:} Dowd et al., Chapters 2--4 (design, operations,
review), Bishop Chapters 19--22 and 24--25 (assurance, vulnerabilities,
auditing), Chess \& West Chapters 3--4 (analysis integration and
reporting).

\textbf{Modern extension:} \url{https://owasp.org/projects/asvs};
\url{https://csrc.nist.gov/pubs/sp/800/218/final};
\url{https://docs.github.com/en/code-security/concepts/code-scanning/code-scanning-alerts};
\url{https://docs.github.com/en/code-security/how-tos/find-and-fix-code-vulnerabilities/configure-code-scanning/configuring-advanced-setup-for-code-scanning};
\url{https://docs.github.com/en/code-security/how-tos/find-and-fix-code-vulnerabilities/integrate-with-existing-tools/upload-sarif-file}.

\end{document}
